%% kapa: exam=1 round=19 year=2008 period=2 form=A subject=economics subject_ko=경제원론 seq=30 number=30
%% source: Q-Net 감정평가사 기출문제 (cst00309) artlSeq=1001962 file=extracted/1008/1001962/2008년 제19회 감정평가사 1차시험 기출문제/2교시.hwp
%% conversion: kordoc 4.24.51 -> tools/gen1.py (auto); review=proofread-sonnet; flags=none
\begin{question}{30}
소국 개방경제(small open economy)인 A국이 해외로부터의 수입 농산물에 관세를 부과할 때 나타나는 현상으로 옳은 것을 모두 고르면? (단, 국내산 농산물과 수입 농산물은 동질적이다. 또한 농산물에 대한 A국의 국내 수요곡선은 우하향하고, 국내 공급곡선은 우상향한다고 가정한다.)
\begin{framedtext}
ㄱ. A국에서 국내 생산에 의한 농산물 공급량이 감소한다. \par
ㄴ. A국에서 해외로부터의 농산물 수입량이 감소한다. \par
ㄷ. A국에서 생산자 잉여는 증가하고 소비자 잉여는 감소한다. \par
\end{framedtext}
\begin{choices}
  \item ㄱ
  \item ㄴ
  \item ㄷ
  \item ㄱ, ㄴ
  \item ㄴ, ㄷ
\end{choices}
\end{question}
\answer{5}
